In this section, we recall and rephrase some results from synthetic Stone duality (SSD, \cite{synthetic-stone-duality}) and synthetic topology \cite{SyntheticTopologyEscardo,SyntheticTopologyLesnik,abstractstone,faissole2017synthetic,vickers2007locales} for the purposes of this note.  

List of stuff we use, to complete later
\begin{itemize}
  \item Open and Closed propositions and subsets. 
  \item Countably presented algebras and some of their closure properties. 
  \item Compact Hausdorff spaces. (Not sure we need ODisc in this note). 
  \item All the axioms (Local choice phrased in terms of CHaus). 
  \item Definition of a basis ala Lesnik.
  \item Point out what is a basis of CHaus spaces. 
  \item Local compactness of CHaus
  \item Product topology. 
  \item Interval, Cauchy completeness, LLPO and MP as interval results. 
  \item Intermediate value theorem. 
\end{itemize}
\subsection{SSD}
\begin{definition}
\item A proposition $P$ is called \textbf{open} if there exists some sequence $\alpha:2^\mathbb N$ such that $P\leftrightarrow \exists_{n:\N} \alpha n = 1$. \parencite[Dfn 1.19]{synthetic-stone-duality}
    \item A proposition $P$ is called \textbf{closed} if there exists some sequence $\alpha:2^\mathbb N$ such that $P\leftrightarrow \forall_{n:\N} \alpha n = 0$. \parencite[Dfn 1.19]{synthetic-stone-duality}
    \item We denote $\Open, \Closed$ for the types of open and closed propositions respectively. \parencite[Dfn 1.19]{synthetic-stone-duality}
  \item For $X$ a set, an \textbf{open subset} of $X$ is a map $X \to \Open$. \parencite[Dfn 1.25]{synthetic-stone-duality}
    \item For $X$ a set, a \textbf{closed subset} of $X$ is a map $X \to \Closed$. \parencite[Dfn 1.25]{synthetic-stone-duality}
    \item A type is \textbf{countable} iff it is merely equivalent to a decidable subset of $\mathbb N$. \parencite[Dfn 1.2]{synthetic-stone-duality}.
    \item A type is \textbf{overtly discrete} if there is the quotient of a countable set by an open equivalence relation. \parencite[Dfn 2.1, Lem 2.12]{synthetic-stone-duality}.
\end{definition}

\subsection{Topology}

\begin{definition}
\item A map $f: A \to B$ is overt if for all $U : B \to \Open$, the statement $\exists_{a:A}.U(f(a))$ is open \parencite[Dfn 2.39]{lesnik2021synthetictopologyconstructivemetric}.
\item A map $f:A \to B$ is compact if for all $U : B \to \Open$, the statement $\forall_{a:A} U(f(a))$ is open \parencite[Dfn 2.39]{lesnik2021synthetictopologyconstructivemetric}. 
\item A subset is subovert/subcompact iff the inclusion map is. 
\item Let $X$ be a set. A \textbf{weak basis} for $X$ is a map $b: \mathcal B\to (X \to \Open)$ such that for any open subset $U\subseteq X$ there merely exists a set $I$ and map $u:I \to \mathcal B$ such that $U = \bigcup_{i: I} b(u_i)$. \parencite[Dfn 2.55]{lesnik2021synthetictopologyconstructivemetric}. 
\item Such a weak basis is called a \textbf{basis} if $u$ can always be taken to be an overt map. 
\item Let $X$ be a set. If there exists a basis of $X$ such that $\mathcal B$ is overtly discrete,
  % and $I$ can always be taken to be $\mathbb N$,
  we call $X$ a \textbf{second countable} space.
\end{definition}
\begin{remark}
  If $A$ is overt or compact, so are all maps out of $A$. In particular, and closed subset of a compact Hausdorff space is subcompact. 
\end{remark}
\begin{lemma}
  Let $X:\Chaus$ with presentation $q:S \twoheadrightarrow X$. 
  Then the map $b:2^S \to (X \to \Open)$ given by $b(D) = \neg q(D)$ forms a basis of $X$. 
  For this basis, one can always take $I= \mathbb N$. 
\end{lemma}
\begin{proof}
  By \parencite[Cor 4.7]{synthetic-stone-duality} and the observation that $\mathbb N$ is overt, and any map from an overt space is overt. 
\end{proof}
\begin{definition}
  Let $X$ be a type. $X$ is \textbf{locally compact} if for any $U\subseteq X$ open and $x\in U$, there exists an open $V\subseteq X$ and subcompact $K\subseteq X$ such that $x \in V\subseteq K \subseteq U$. \parencite[p59]{lesnik2021synthetictopologyconstructivemetric}
\end{definition}  
\begin{lemma}
  If $X$ is locally compact, then $X\times Y$ has a weak basis given by 
  $(X \to \Open) \times (Y \to \Open)\to (X \times Y) \to \Open$ sending $U\subseteq X , V \subseteq Y$ to $U\times V$. 
  \parencite[Prop 2.58]{lesnik2021synthetictopologyconstructivemetric}.
\end{lemma}
\begin{corollary}
  If $X,Y$ both have a basis and $X$ is locally compact, then 
\end{corollary}

\subsection{Analysis}




\begin{remark}\label{IntermediateValue}
  The intermediate value theorem as stated in Theorem 5.29 of \cite{synthetic-stone-duality} has the following consequence: If $f : [0,1] \to [0,1]$ is such that for some $x<y, a<b$ we have $f(x) = a, f(y) = b$, then for all $c$ with $a<c<b$ there exists some $z$ with $x<z<y$ such that $f(z) = c$. 
\end{remark}
\begin{remark}
  Closed propositions are double negation stable and closed under disjunction and countable intersection. 
\end{remark}
\begin{remark}
  Fix the statement of $\Sigma(x \sim y) $ vs $\exists(x\sim y)$ for $x\sim y$ the path relation. 
\end{remark}
\begin{remark}
  Cauchy reals agree with set quotient of $2^\mathbb N$ by $\sim_n$. 
  Conclude from this which functions one can definitely define on $\mathbb R^n \to \mathbb R^m$. 
\end{remark}

\begin{remark}
  Let $X:\CHaus$ be represented by $q:S\twoheadrightarrow X$. Lemma \parencite[TODO]{synthetic-stone-duality} says that any open $U\subseteq X$ can be written as $\bigcup_{n:\mathbb N} \neg q(D_n)$ where $D : \mathbb N \to 2^S$ is some $\mathbb N$-indexed collection of decidable subsets of $S$. 
\end{remark}
